\clearpage
\VerdierPhysicalPageAnchor{0171}%
\begingroup\small\setlength{\parskip}{0pt}\setlength{\emergencystretch}{1.5em}
\VerdierSemanticAnchor{ch4-4-3-3}
\medskip\noindent\textbf{4.3.3.}\quad
特に項
$E\!\left[\begin{smallmatrix}n\\p-r+1,p,p+1,p+r\end{smallmatrix}\right]$、
$p\in\mathbb Z$、$r\in\mathbb Z$、$r>0$ に注目する。適切な仮定のもとで、
これらから $H^n(-\infty,+\infty)$ に付随する次数付き対象を決定できるからで
ある。\emph{同所引用文献}の記法に戻すため、次のように置く：

\VerdierSemanticAnchor{ch4-4-3-3-1}
\[
\left\{
\begin{aligned}
E\!\left[\begin{smallmatrix}n\\p-r,p-1,p,p+r-1\end{smallmatrix}\right]
  &=\mathrm I_r^{-p,n+p},\\
&\textit{すなわち}:\\[-2pt]
\mathrm I_r^{p,q}
  &=E\!\left[\begin{smallmatrix}p+q\\-p-r,-p-1,-p,-p+r-1\end{smallmatrix}\right],
  &&p,q\in\mathbb Z,\quad 1\leq r\leq+\infty,\\
&\text{対応する微分を次のように記す：}\\[-2pt]
d_r\colon\mathrm I_r^{p,q}&\longrightarrow
  \mathrm I_r^{p+r,q-r+1}.
\end{aligned}
\right.
\tag{4.3.3.1}
\]

次の記法も非常によく用いる：

\VerdierSemanticAnchor{ch4-4-3-3-2}
\[
\left\{
\begin{aligned}
E\!\left[\begin{smallmatrix}n\\p-r+1,p,p+1,p+r\end{smallmatrix}\right]
  &=\mathrm{II}_{r+1}^{n-p,p},\\
&\textit{すなわち}:\\[-2pt]
\mathrm{II}_r^{p,q}
  &=E\!\left[\begin{smallmatrix}p+q\\q-r+2,q,q+1,q+r-1\end{smallmatrix}\right],
  &&p,q\in\mathbb Z,\quad 2\leq r\leq+\infty,\\
&\text{対応する微分を次のように記す：}\\[-2pt]
d_r\colon\mathrm{II}_r^{p,q}&\longrightarrow
  \mathrm{II}_r^{p+r,q-r+1}.
\end{aligned}
\right.
\tag{4.3.3.2}
\]

各場合に、どちらの記法を用いるかを明記する。記法 (4.3.3.1) を用いるとき、
項 $\mathrm I_r^{p,q}$ は $r\geq1$ に対して定義されるが、記法
(4.3.3.2) を用いるとき、項 $\mathrm{II}_r^{p,q}$ は $r\geq2$ に対して
のみ定義されることに注意する。

\VerdierSemanticAnchor{ch4-4-3-3-spectral-sequence}
項
$E\!\left[\begin{smallmatrix}n\\p-r+1,p,p+1,p+r\end{smallmatrix}\right]$、
$n,p\in\mathbb Z$、$1\leq r\leq+\infty$ の全体に対応する微分を備えたものを、
スペクトル対象 $(H^n(p,q),\delta(p,q,r))$ に\emph{付随するスペクトル系列}
と呼ぶ。
\endgroup
